\documentclass{article}
\usepackage{iclr2027_conference,times}
\usepackage{amsmath,amssymb}
\usepackage[colorlinks=true,citecolor=black,linkcolor=black,urlcolor=blue]{hyperref}

\iclrfinalcopy

\title{Exploration-Aware Latent World Models for MuZero Planning\\[0.4em]
\large Thesis Research Proposal}
\author{Pawel Swiderski}
\date{\today}

\begin{document}
\maketitle
\lhead{Thesis Research Proposal}

\begin{abstract}
Sparse-reward tasks may require coordinated action sequences before any task
reward is observed. MuZero can guide tree search using learned intrinsic reward
and value predictions, but the RND signal does not decrease when an observation
is revisited before the next network update. Outcomes not yet experienced also
provide no direct reward supervision for the learned intrinsic heads.
This thesis investigates a separate observation-grounded world model that
predicts the observations resulting from candidate action sequences. Predicted
observation embeddings are compared with real episodic memory and temporary
branch memory to estimate whether an outcome has already been encountered.
Predicted observations are also evaluated by the current Random Network
Distillation (RND) networks, providing an alternative estimate of future
intrinsic reward. The study will compare this estimate with MuZero's learned
intrinsic reward head and test combinations selected on development layouts.
Experiments in partially observed, sparse-reward navigation will evaluate
revisit detection, RND-reward estimation, coverage, and hidden-goal discovery
under matched interaction and search budgets. The analysis will distinguish
improvements in prediction accuracy from improvements in exploration.
\end{abstract}

\section{Introduction and Motivation}

Sparse-reward problems may require an agent to commit to a sequence of actions
before receiving useful feedback from the task. Effective exploration can
therefore require passing through familiar states to reach unexplored regions,
which choosing random actions independently may fail to accomplish
\citep{bootstrapped-dqn}.

MuZero is a natural starting point because its MCTS evaluates possible action
sequences in a learned latent model \citep{muzero}. An intrinsic reward can be
added to the task reward, with separate reward and value heads used to predict
its effect during search. Random Network Distillation (RND), for example,
gives a high intrinsic reward when a learned predictor has high error on an
observation \citep{rnd}. After the agent has received such a reward, the
intrinsic value head can learn the returns from preceding states, allowing that
reward to influence planning even when its location lies beyond the explicit
search horizon. \citet{walker23} provide a direct precedent for training a
MuZero-style model with RND and distinguish the intrinsic signal predicted
during planning from the signal measured on real observations.

During data collection, the RND predictor and MuZero networks are held fixed
for a sequence of actions. Encountering the same observation twice before the
next update therefore produces exactly the same RND reward both times.
MuZero's intrinsic heads are also unchanged between the two visits, so they
cannot learn from the first visit before scoring the second. Updating RND alone
would still leave the predictions used by MuZero during search unchanged.
Furthermore, outcomes not yet experienced provide no direct reward supervision
for MuZero's intrinsic heads, although the heads may generalize from related
observations.

This thesis investigates a separate observation-grounded world model that
predicts the observations resulting from candidate action sequences. Its
predicted embeddings are compared with episodic memory to estimate revisits,
while its decoded observations are evaluated by the current RND networks to
estimate future intrinsic rewards. This direct estimate does not require a
previously collected RND reward label for the candidate outcome, but it depends
on accurate observation prediction. The study will compare direct and
head-based RND estimates across planning depths and evaluate combinations
selected on development layouts.

Intrinsic rewards already shape MuZero's latent state: once intrinsic reward
and value are training targets, the latent must retain the information needed
to predict them. The problem is narrower. Predicting scalar rewards and values
does not require distances in the latent space to tell us whether two states
are the same. Different locations can have similar intrinsic returns, and the
same location is novel early in an episode but familiar after a visit. The
separate model is introduced to support observation comparison and direct
evaluation of imagined outcomes. MuZero's intrinsic value head keeps its
existing role of carrying experienced RND rewards beyond the explicit search
horizon.

\section{Related Work}

MuZero learns representation, dynamics, reward, policy, and value functions
for tree search without reconstructing the environment state \citep{muzero}.
The value-equivalence principle explains why a model can preserve the
information needed for its value predictions while discarding other state
details \citep{value-equivalence}. This can be useful for planning, but scalar
task and intrinsic targets do not require distances in the latent space to
identify repeated states.

RND uses prediction error against fixed random features as a measure of
long-term novelty \citep{rnd}. Never Give Up (NGU) adds an episodic memory so
that a state becomes less rewarding after it has been visited in the current
episode. Agent57 uses the same general mechanism in a multi-policy agent
\citep{ngu,agent57}. These methods calculate episodic novelty after receiving
the next real observation. In MCTS, the agent must instead estimate the
familiarity of a state that has only been imagined.

\citet{walker23} combine a MuZero-style model with RND for intrinsic
exploration, using learned predictions of the intrinsic reward during
planning. The present study uses this head-based estimate as a baseline
and asks whether applying the current RND networks directly to a predicted
observation provides a more accurate edge reward.

World models such as PlaNet and Dreamer learn recurrent latent dynamics from
visual trajectories and use predicted trajectories for control
\citep{planet,dreamer}. Plan2Explore uses a learned world model to seek outcomes
that are expected to be informative \citep{plan2explore}. The model proposed
here has a more specific role alongside MuZero: it supplies predicted
observation representations for episodic comparison and predicted observations
on which the current RND networks can be evaluated. MuZero continues to provide
task predictions and the learned intrinsic value function.

GoBI also combines a learned world model with episodic memory, rewarding
expansion of the set of real and predicted reachable states \citep{gobi}. In the
present approach, predicted observations are used only while evaluating an
MCTS branch. They enter temporary branch memory but are never committed to the
episode's persistent memory.

\section{Problem Formulation}

\subsection{Task and intrinsic prediction in MuZero}

Let $o_t\in\mathcal{O}$ be the observation at time $t$, let
$a_t\in\mathcal{A}$ be the selected action, and let $H_t$ denote the
observation--action history through $o_t$. MuZero represents the history and
recurrently predicts imagined transitions:
\begin{align}
    s_t &= h_\theta(H_t), \\
    (\hat{s}_{t+1},\hat{r}^{\mathrm{task}}_{t+1},
      \hat{r}^{\mathrm{int}}_{t+1})
        &= g_\theta(s_t,a_t), \\
    (\hat{p}_t,\hat{v}^{\mathrm{task}}_t,
      \hat{v}^{\mathrm{int}}_t)
        &= f_\theta(s_t).
\end{align}
Here $s_t$ is the MuZero latent, $\hat{p}_t$ is the policy prior, and the
superscripts distinguish task and intrinsic predictions. The heads may share
some network layers, but the distinction that matters here is between their
training targets: the intrinsic heads learn from intrinsic rewards collected
through real interaction.

For RND, a fixed target network $f^\ast$ and a learned predictor $f_\eta$
define the intrinsic reward of the next observation:
\begin{equation}
    r^{\mathrm{RND}}_{t+1}
      = \left\|f_\eta(o_{t+1})-f^\ast(o_{t+1})\right\|_2^2.
\end{equation}
The intrinsic reward head predicts this quantity. The intrinsic value head
predicts the discounted sum of future intrinsic rewards. For a simulated
branch of depth $L$, the intrinsic part of the MCTS return is
\begin{equation}
    \hat{G}^{\mathrm{int}}_t
      = \sum_{k=1}^{L}\gamma^{k-1}
        \hat{r}^{\mathrm{int}}_{t+k}
        +\gamma^L\hat{v}^{\mathrm{int}}_{t+L},
\end{equation}
where $\gamma$ is the discount factor. This value term is what allows an
intrinsic reward discovered at a distant state to affect earlier decisions
after the agent has learned from it.

\subsection{Two missing signals}

Let $c$ identify one collection phase. The parameters $\eta_c$ and $\theta_c$
are the RND and MuZero parameters used during that phase. The agent collects a
sequence of transitions with these parameters and updates the networks
afterwards. If it encounters the same observation at times $i$ and $j$ before
the update, the RND reward is unchanged:
\begin{equation}
    o_i=o_j
    \quad\Longrightarrow\quad
    r^{\mathrm{RND}}_{\eta_c}(o_i)
      =r^{\mathrm{RND}}_{\eta_c}(o_j).
\end{equation}
Updating $\eta_c$ immediately would change the reward on the real observation,
but it would not update the reward and value predictions produced by
$\theta_c$. Synchronizing both components after every visit would require
additional RND and MuZero optimization between consecutive actions while
continually moving the intrinsic training targets. RND therefore measures
familiarity learned before the current collection phase. It does not
immediately record that an observation has just been encountered in the
current episode.

The reward head will not be trained on historical RND labels stored at the time
of collection. After each collection phase, the RND predictor is updated on the
new observations. A frozen copy of the updated predictor is then used together
with the fixed target network $f^\ast$ to recompute RND reward and value targets
from replay observations whenever a batch is sampled. The same frozen networks
define direct RND evaluation and the realized target used to compare the two
estimators. Their comparison therefore measures differences between the
estimators rather than differences in the age of their reward labels.

The second problem concerns an outcome that has never been observed. For a
candidate branch ending in $o_{t+k}$, neither
$r^{\mathrm{RND}}(o_{t+k})$ nor its preceding intrinsic return is present in
the replay data before the first visit. The learned heads can extrapolate from
related experience, but direct supervision for that particular outcome becomes
available only after it is observed. An observation-grounded model offers a
different route: predict the future observation and apply the current RND
target and predictor to that prediction. This does not require a previously
observed reward label for the candidate outcome, but it introduces error from
the world model.

Let $x_t=E_\phi(o_t)$ be the exploration embedding of the individual real
observation $o_t$. The episodic memory contains only embeddings of observations
received during the current episode:
\begin{equation}
    \mathcal{M}^{\mathrm{ep}}_t
        =\{\operatorname{norm}(x_i):
          i\leq t\text{ in the current episode}\}.
\end{equation}
It is reset at the episode boundary. Given a predicted future observation
embedding $\tilde x_{t,k}$, this memory asks whether a similar observation has
already occurred in the episode. It does not claim to measure lifetime novelty
or epistemic uncertainty.

\section{Research Question and Hypotheses}

\textbf{Research question.} Can a separate observation-grounded world model
improve exploration in MuZero by identifying imagined episodic revisits and
estimating future RND rewards through direct evaluation of predicted
observations?

\textbf{Representation hypothesis.} Predicted observation embeddings will
identify same-episode physical-pose revisits more accurately than comparisons
made with MuZero's latent states.

\textbf{Intrinsic-reward hypothesis.} At planning depths where future
observations can be predicted accurately, evaluating the current RND networks
on those predictions will estimate the RND reward of the actual future
observation more accurately than the learned intrinsic reward head, especially
for outcomes not represented in its reward supervision.

\textbf{Behavioral hypothesis.} Reducing the intrinsic reward of an imagined
episodic revisit will reduce repeated local movement. Replacing or correcting
the learned edge-reward estimate with a more accurate direct RND estimate will
improve the selection of exploratory action sequences. Together with the
existing intrinsic value, these changes will improve coverage and sparse-goal
discovery under a fixed interaction budget.

These claims will be tested separately: episodic correction may help even when
direct RND does not outperform the reward head, and more accurate predictions
need not improve action selection.

\section{Proposed Approach}

\subsection{Observation-grounded exploration model}

The exploration model contains an observation encoder $E_\phi$, recurrent
belief updater $F_\phi$, decoder $D_\phi$, and action-conditioned transition
model $T_\phi$. A projection $P_\phi$ maps the recurrent prediction back into
the observation-embedding coordinates. For a real trajectory,
\begin{align}
    x_t &= E_\phi(o_t), \\
    b_t &= F_\phi(b_{t-1},x_t,a_{t-1}), \\
    \hat{o}_t &= D_\phi(x_t), \\
    z_t &= \operatorname{norm}(x_t).
\end{align}
The recurrent belief $b_t$ incorporates observation and action history for
prediction. The vector $z_t$ is the normalized observation embedding stored in
episodic memory and used for nearest-neighbour comparison. It represents the
individual observation rather than the complete history. This follows the
practical separation used by episodic-memory agents: history helps predict what
comes next, while memory records what the agent actually observed
\citep{ngu}.

Starting from the real root belief $b_t$, a candidate action sequence is
predicted recursively without access to future observations:
\begin{align}
    \tilde b_{t,0} &= b_t, \\
    \tilde b_{t,k}
        &=T_\phi(\tilde b_{t,k-1},a_{t+k-1}), \\
    \tilde x_{t,k} &=P_\phi(\tilde b_{t,k}), \\
    \tilde o_{t,k} &=D_\phi(\tilde x_{t,k}), \\
    \tilde z_{t,k}
        &=\operatorname{norm}(\tilde x_{t,k}),
        \qquad k=1,\ldots,K.
\end{align}

The model is trained on real replay sequences with three separate objectives:
\begin{align}
    \mathcal{L}_{\mathrm{rec}}
        &=\mathbb{E}\left[
          \ell_{\mathrm{obs}}(D_\phi(x_t),o_t)
          \right], \\
    \mathcal{L}_{\mathrm{latent}}
        &=\mathbb{E}\left[
          \frac{1}{K}\sum_{k=1}^{K}
          \left\|\tilde z_{t,k}
          -\operatorname{sg}(z_{t+k})\right\|_2^2
          \right], \\
    \mathcal{L}_{\mathrm{future}}
        &=\mathbb{E}\left[
          \frac{1}{K}\sum_{k=1}^{K}
          \ell_{\mathrm{obs}}(\tilde o_{t,k},o_{t+k})
          \right], \\
    \mathcal{L}_{\mathrm{expl}}
        &=\lambda_{\mathrm{rec}}\mathcal{L}_{\mathrm{rec}}
          +\lambda_{\mathrm{latent}}\mathcal{L}_{\mathrm{latent}}
          +\lambda_{\mathrm{future}}\mathcal{L}_{\mathrm{future}}.
\end{align}
The reconstruction loss encourages the encoder to preserve observation
information. The latent and future-observation losses train recursive
predictions to match the future embedding and image, respectively. The image
prediction loss is needed because direct RND evaluation operates on decoded
observations. The stop-gradient operator blocks gradients through the target
embedding in the latent-prediction loss. The representation continues to learn
between episodes. During an episode, the exploration model is held fixed, so
all entries in episodic memory use the same encoder. The memory is reset at the
episode boundary. Physical coordinates and room identity are not used for
training.

\subsection{Memory-conditioned novelty}

Let $N(z,\mathcal{M})\in[0,1]$ be a normalized novelty score that increases as
$z$ becomes less similar to the nearby entries in memory $\mathcal{M}$. The
primary implementation will use a density estimate based on nearest-neighbour
distances, with its scale estimated from distances among real memory entries
\citep{ngu}.
The exact density function is secondary. For the method to work, repeated or
similar observations must receive a low score, while observations far from the
stored examples must receive a high score.

When MCTS evaluates a node at depth $k$, let
\begin{equation}
    \mathcal{B}_{t,k}
      =\{\tilde z_{t,1},\ldots,\tilde z_{t,k-1}\}
\end{equation}
contain the observation embeddings predicted earlier on that same branch.
Episodic novelty is
\begin{equation}
    n^{\mathrm{ep}}_{t,k}
        =N\left(
          \tilde z_{t,k},
          \mathcal{M}^{\mathrm{ep}}_t\cup\mathcal{B}_{t,k}
          \right).
\end{equation}
Thus, the predicted observation at depth $k$ is compared with both real
episodic memory and the earlier predictions on its own branch. Suppose a branch
predicts the sequence $A,B,A$. The first predicted visit to $A$ may receive a
novelty bonus,
but the final visit will be compared with the earlier predicted $A$ and should
receive a low score. Without this temporary branch memory, an imagined loop
through states absent from real memory could receive the novelty bonus on every
visit. The set $\mathcal{B}_{t,k}$ is used only while evaluating that branch
and is discarded after search. Only embeddings of real observations are kept
in episodic memory.

\subsection{Two estimates of future RND reward}

The learned intrinsic reward head and direct evaluation of the predicted
observation estimate the same target. For a future observation that is revealed
by a logged real trajectory, define that target using the RND parameters
available at the search root:
\begin{equation}
    r^{\mathrm{actual}}_{t,k}
      =\left\|f_\eta(o_{t+k})-f^\ast(o_{t+k})\right\|_2^2.
\end{equation}
The ordinary head-based estimate and the direct model-based estimate are
\begin{align}
    \hat r^{\mathrm{head}}_{t,k}
      &=\hat r^{\mathrm{int}}_{t+k}, \\
    \hat r^{\mathrm{direct}}_{t,k}
      &=\left\|f_\eta(\tilde o_{t,k})
                    -f^\ast(\tilde o_{t,k})\right\|_2^2.
\end{align}
The first estimate is learned from RND rewards on real transitions in replay.
The second composes the exploration model with the current RND networks. It can
score a predicted outcome even when the corresponding real future observation
has never supplied an RND reward label in replay. However, its score is useful
only if $\tilde o_{t,k}$ resembles the observation that the environment would
actually produce. Prediction artefacts may receive high RND error even when the
real outcome is familiar, while a prediction that looks too familiar may hide
the novelty of the real outcome.

Because these are estimates of one target, they are compared rather than added.
Both are first expressed in the raw RND-reward scale by inverting any target
transform used to train the reward head. A family of edge estimates is then
defined by
\begin{equation}
    \tilde r^{\mathrm{int}}_{t,k}(w_k)
      =n^{\mathrm{ep}}_{t,k}
       \left[(1-w_k)\hat r^{\mathrm{head}}_{t,k}
             +w_k\hat r^{\mathrm{direct}}_{t,k}\right],
    \qquad 0\leq w_k\leq1.
\end{equation}
The setting $w_k=0$ uses the reward head, $w_k=1$ uses direct RND evaluation,
and intermediate values define hybrids. Depth-dependent weights may be selected
on development layouts because direct observation prediction is expected to
become less reliable deeper in the tree. Development diagnostics will compare
the head-based and direct estimates separately before hybrid weights are
selected. Final diagnostics will report both estimates at the prespecified
depths.

For a simulated branch of depth $L$, MCTS backs up
\begin{equation}
    \hat G_t
      =\sum_{k=1}^{L}\gamma^{k-1}
       \left(
       \hat r^{\mathrm{task}}_{t+k}
       +\beta\tilde r^{\mathrm{int}}_{t,k}(w_k)
       \right)
       +\gamma^L\left(
       \hat v^{\mathrm{task}}_{t+L}
       +\beta\hat v^{\mathrm{int}}_{t+L}
       \right).
\end{equation}
The intrinsic value head remains trained on ungated RND returns. Episodic
correction therefore applies only to the explicitly simulated part of each
branch; the leaf value may still assign value to outcomes already visited in
the current episode. It nevertheless carries experienced intrinsic reward
beyond the depth of the search tree. Direct RND evaluation is also limited to
explicitly predicted outcomes and cannot reveal an unseen region beyond the
exploration model's useful prediction horizon. Separate ablations remove
episodic gating, use each edge estimator alone, and test the selected hybrid.

\section{Main Experiment}

\subsection{Branching Rooms benchmark}

The main experiment will use a family of Branching Rooms environments built
with
MiniGrid \citep{minigrid}. Each layout contains visually marked rooms connected
by narrow corridors, with branches, loops, bottlenecks, and dead ends. The
agent receives a local RGB observation and can turn left, turn right, or move
forward. It does not observe its coordinates or the global room graph.
Transitions are deterministic. Because the agent sees only part of the maze
and different places may look alike, it must use its history to distinguish
some states.

Each episode begins in a designated root room. An unmarked terminal goal is
placed in a distant dead end; reaching it yields task reward $1$, and every
other transition yields zero task reward. Before the goal is found, the task
reward itself supplies no evidence about which exploratory route is useful.
Necessary corridors and occasional backtracking also ensure that simply
penalizing every revisit is not an adequate solution.

Layouts will be generated from fixed seeds and divided into development and
held-out test sets. Development layouts are used to choose the model design,
memory scale, intrinsic-reward weights, prediction depths, and other
hyperparameters. Once these choices are fixed, they are not changed in response
to results on the test layouts. Learned weights and replay data from the
development layouts are not transferred to the final trials.

On each held-out layout, every method is initialized from scratch and trained
for $T$ environment interactions. MuZero, the exploration model, the RND
target and predictor, optimizer state, replay, recurrent state, and episodic
memory are all initialized anew for each seed--layout trial. Methods are paired
by layout and random seed, and common components use matched initializations
where possible.
Nothing learned in one test trial carries into another. The resulting learning
curves therefore measure exploration and learning within a previously unseen
layout, rather than transfer from a pretrained checkpoint. Every method
receives the same number of interactions and MCTS simulations. Runtime and
action-selection latency will also be reported because the additional model
increases computation.

\subsection{Predictive evaluation}

The final runs will include prespecified predictive diagnostics that are not
used to change the method or its hyperparameters. On evaluation-only logged
trajectories, episodic familiarity and the two RND estimators will be tested at
fixed training checkpoints. At a selected search root $t$, the model receives
the history up to that point and the next logged actions, but not the future
observations. It predicts $\tilde z_{t,k}$ and $\tilde o_{t,k}$ for depths $1$
through $K$. The logged trajectory later reveals the observation and physical
pose actually reached at each depth.

For the primary revisit label, the physical pose consists of grid position and
orientation. The future pose is a revisit if it appeared at or before the
selected root or at an earlier step of the logged action suffix. Position-only
and room coverage are reported separately. Because
$n^{\mathrm{ep}}_{t,k}$ is high for an outcome predicted to be new, the revisit
score is $1-n^{\mathrm{ep}}_{t,k}$. Its area under the receiver operating
characteristic curve (AUC) is reported at every depth. An AUC of $1$ means that
every revisit is ranked above every non-revisit; an AUC of $0.5$ is random
ranking.

The RND test compares $\hat r^{\mathrm{head}}_{t,k}$ and
$\hat r^{\mathrm{direct}}_{t,k}$ with
$r^{\mathrm{actual}}_{t,k}$ at every depth. Both estimators and the realized
target use the same RND predictor and target networks, held fixed at the
diagnostic checkpoint. The analysis reports mean absolute error, Spearman rank
correlation, and plots of mean predicted against realized reward within
prediction bins. Results are stratified in two ways: whether the physical pose
has previously been visited in the current test trial, and whether the exact
observation has previously appeared in replay available for reward-head
training. The second distinction identifies outcomes for which no direct reward
label is available. Simulator pose identifiers are used only to create
evaluation labels and are not given to the agent or used for training.

The episodic test is run in three ways:
\begin{itemize}
    \item using the embedding of the true future observation, which jointly
    tests the observation representation, memory lookup, and perceptual
    ambiguity in the observation itself;
    \item using the predicted future observation embedding, which also tests
    the transition model; and
    \item using MuZero's latent state, which tests whether a separate
    representation is needed.
\end{itemize}
In the MuZero-latent control, each real memory entry is the normalized root
state $h_\theta(H_i)$ produced from the observation history at time $i$. During
MCTS, the query is the normalized recurrent state produced by applying
$g_\theta$ along the candidate action sequence. The episodic memory, temporary
branch memory, novelty function, and use of the episodic factor are otherwise
unchanged. This control therefore replaces only the source of the stored and
predicted embeddings.

The target $r^{\mathrm{actual}}_{t,k}$ is the current RND score of the true
future observation and is available only for offline evaluation. Comparing it
with direct RND on the predicted observation measures the error introduced by
the exploration model.

\subsection{Behavioral evaluation}

The core behavioral study crosses two edge estimators, head-based and direct,
with episodic gating either disabled or enabled. The development-selected
hybrid is also evaluated with both gating settings. The comparisons therefore
include:
\begin{itemize}
    \item the learned reward head, with and without episodic gating;
    \item direct RND evaluation, with and without episodic gating;
    \item the selected hybrid, with and without episodic gating; and
    \item a head-based control that applies episodic gating using MuZero's
    latent state instead of the separate observation embeddings.
\end{itemize}

At fixed checkpoints during each held-out-layout training trial, a frozen copy
of the current agent will be evaluated for a fixed number of episodes.
Evaluation episodes do not enter replay and do not update MuZero, the
exploration model, or the RND predictor. Let
$0=e_0<\cdots<e_Q=T$ be the evaluation checkpoints, where $T$ is the total
training-interaction budget on that layout, and let $S_q\in[0,1]$ be the
fraction of evaluation episodes in which the hidden goal is reached at
checkpoint $q$. The primary
behavioral measure is the normalized area under this learning curve:
\begin{equation}
    A_{\mathrm{success}}
      =\frac{1}{T}\sum_{q=0}^{Q-1}
       \frac{S_q+S_{q+1}}{2}
       \left(e_{q+1}-e_q\right).
\end{equation}
This is the average success rate over the course of training. An agent that
learns early receives a higher score than one that reaches the same final
success rate only near the end. This learning-curve score is different from
the receiver operating characteristic AUC used in the predictive tests.
Secondary measures are the time to first goal, final success rate, position and
room coverage, distant dead ends reached, and the fraction of transitions that
return to a physical pose already visited in the same episode. All methods use
the same random seeds and test layouts.

The episodic correction should mainly reduce repeated local movement. Direct
RND evaluation should help only if the exploration model predicts enough of the
future observation for its RND score to be more accurate than the learned
head. The selected method must improve goal discovery, not just its own
prediction metric. A method that avoids necessary backtracking, for example,
may revisit fewer poses but perform worse on the task.

\section{Decision Rule}

The representation hypothesis is supported if predicted observation embeddings
identify episodic revisits more accurately than MuZero's latent states at the
planning depths chosen before final evaluation. Results obtained from true
future observations and predicted embeddings will be reported separately. This
separates limitations in observation-based lookup from errors introduced by
the transition model.

The episodic mechanism is supported if adding only episodic correction reduces
redundant same-episode revisits without reducing distant coverage or goal
discovery. Direct RND evaluation is supported if it estimates the actual future
RND reward more accurately than the intrinsic reward head at prespecified
planning depths. A behavioral benefit from direct estimation requires the
direct or hybrid estimator to improve the hidden-goal learning curve over the
reward head under the same episodic-gating setting. The complete method will
also be compared with the original ungated, head-based MuZero--RND agent under
the same interaction and search budgets.

The main comparisons will pair methods by random seed and layout. They will
report effect sizes and confidence intervals, together with a paired
randomization or permutation test for the chosen outcomes. If the model
predicts familiarity accurately but does not improve behavior, familiarity may
not be the limiting factor, or the signal may need to be integrated differently
into MCTS. If behavior improves without a corresponding improvement in the
relevant prediction, the result does not support the proposed mechanism. If
only the episodic correction works, the conclusion will be limited to
controlling revisits within an episode. If direct RND is more accurate but does
not improve behavior, the result supports the estimator but not its use for
action selection.

\section{Expected Contribution and Risks}

The thesis will establish when observation-grounded predictions provide useful
exploration signals for MuZero, distinguishing improvements in episodic revisit
detection, future RND estimation, and exploratory behavior.

The main risks are:
\begin{itemize}
    \item \textbf{Prediction outside experience.} The world model may be least
    accurate exactly where direct RND evaluation is intended to help.
    Comparing RND on true and predicted future observations will reveal this
    problem. Development results will determine the hybrid's depth-dependent
    weights, which will remain fixed during final trials. The direct-only
    variant will remain a separate experimental condition.

    \item \textbf{Observation similarity may not equal a revisit.} Two
    physical poses can produce the same image, while irrelevant visual
    differences can make one pose produce different embeddings. Either
    error can make episodic gating suppress the wrong action. Revisit labels
    based on physical pose will measure how often this occurs.

    \item \textbf{Necessary revisits.} Corridors and backtracking make some
    revisits necessary. Even though episodic gating only reduces an intrinsic
    bonus, it may still make a required return path less attractive and reduce
    goal discovery. Coverage and task performance will therefore be evaluated
    alongside revisit rate.

    \item \textbf{Computational cost.} The additional recurrent model and
    memory queries make planning slower. Results with equal environment
    interactions will therefore be accompanied by latency measurements and a
    second comparison with the same planning-time budget.
\end{itemize}

\section{Provisional Timeline}

\begin{itemize}
    \item \textbf{Months 1--2:} implement and validate the Branching Rooms
    benchmark family and hidden-goal task; verify the existing MuZero--RND
    baseline with a minimal end-to-end run; and implement the exploration
    representation and depth-wise prediction tests.

    \item \textbf{Month 3:} measure episodic-memory accuracy and compare
    head-based with direct future-RND estimates at each prediction depth.

    \item \textbf{Months 4--5:} integrate episodic gating and the validated RND
    estimators with MCTS, then run the paired behavioral comparisons and
    ablations.

    \item \textbf{Month 6:} complete the statistical analysis, document
    failure cases and limitations, and prepare the thesis and reproducibility
    materials.
\end{itemize}

\begin{thebibliography}{99}

\bibitem[Schrittwieser et al.(2020)]{muzero}
J. Schrittwieser et al.
\newblock Mastering Atari, Go, chess and shogi by planning with a learned model.
\newblock \emph{Nature}, 588:604--609, 2020.

\bibitem[Grimm et al.(2020)]{value-equivalence}
C. Grimm, A. Barreto, S. Singh, and D. Silver.
\newblock The value equivalence principle for model-based reinforcement learning.
\newblock In \emph{Advances in Neural Information Processing Systems}, 33, 2020.

\bibitem[Osband et al.(2016)]{bootstrapped-dqn}
I. Osband, C. Blundell, A. Pritzel, and B. Van Roy.
\newblock Deep exploration via bootstrapped DQN.
\newblock In \emph{Advances in Neural Information Processing Systems}, 29, 2016.

\bibitem[Burda et al.(2019)]{rnd}
Y. Burda, H. Edwards, A. Storkey, and O. Klimov.
\newblock Exploration by random network distillation.
\newblock In \emph{International Conference on Learning Representations}, 2019.

\bibitem[Walker et al.(2023)]{walker23}
J. C. Walker, E. V\'{e}rtes, Y. Li, G. Dulac-Arnold, A. Anand,
T. Weber, and J. B. Hamrick.
\newblock Investigating the role of model-based learning in exploration and transfer.
\newblock In \emph{Proceedings of the 40th International Conference on Machine Learning},
pages 35368--35383, 2023.

\bibitem[Puigdom\`enech Badia et al.(2020b)]{ngu}
A. Puigdom\`enech Badia et al.
\newblock Never give up: Learning directed exploration strategies.
\newblock In \emph{International Conference on Learning Representations}, 2020b.

\bibitem[Puigdom\`enech Badia et al.(2020a)]{agent57}
A. Puigdom\`enech Badia et al.
\newblock Agent57: Outperforming the Atari human benchmark.
\newblock In \emph{Proceedings of the 37th International Conference on Machine Learning}, pages 507--517, 2020a.
c
\bibitem[Hafner et al.(2019)]{planet}
D. Hafner et al.
\newblock Learning latent dynamics for planning from pixels.
\newblock In \emph{Proceedings of the 36th International Conference on Machine Learning}, pages 2555--2565, 2019.

\bibitem[Hafner et al.(2020)]{dreamer}
D. Hafner, T. Lillicrap, J. Ba, and M. Norouzi.
\newblock Dream to control: Learning behaviors by latent imagination.
\newblock In \emph{International Conference on Learning Representations}, 2020.

\bibitem[Sekar et al.(2020)]{plan2explore}
R. Sekar et al.
\newblock Planning to explore via self-supervised world models.
\newblock In \emph{Proceedings of the 37th International Conference on Machine Learning}, pages 8583--8592, 2020.

\bibitem[Fu et al.(2023)]{gobi}
Y. Fu, R. Peng, and H. Lee.
\newblock Go beyond imagination: Maximizing episodic reachability with world models.
\newblock In \emph{Proceedings of the 40th International Conference on Machine Learning},
pages 10405--10420, 2023.

\bibitem[Chevalier-Boisvert et al.(2023)]{minigrid}
M. Chevalier-Boisvert et al.
\newblock Minigrid \& Miniworld: Modular \& customizable reinforcement learning environments for goal-oriented tasks.
\newblock In \emph{Advances in Neural Information Processing Systems}, 36, 2023.

\end{thebibliography}

\end{document}
